\chapter{Mechanische energie en het behoud ervan}\label{ch:g11:mechanical-energy}

Een achtbaan heeft voorbij de eerste heuvel geen motor meer: elke lus,
elke uitbarsting van snelheid daarna wordt betaald uit de rekening die op
de weg omhoog is geopend. Dit hoofdstuk brengt die boekhouding terug tot
twee getallen --- kinetische en \omterm{def:g11:mechanical-energy:potential}{potentiële energie} --- en één regel:
zolang alleen het \omterm{def:g10:universal-gravitation:weight}{gewicht} \omterm{def:g11:work-of-force:work}{arbeid} verricht, verandert hun som niet. De
\omterm{def:g10:inertia:inventory}{wrijving} is de kostenpost; we leren die uit de boeken af te lezen.

\section{Kinetische energie}

\begin{definition}[Kinetische energie]\label{def:g11:mechanical-energy:kinetic}
Een lichaam met massa $m$ (\unit{kg}) dat met snelheid $v$ (\unit{m/s})
beweegt, draagt de \emph{kinetische energie}\index{kinetische energie}
\[
E_k = \tfrac12\, m v^2 .
\]
Controle van de eenheden: $\unit{kg.m^2/s^2} = \unit{N.m} = \unit{J}$ ---
\omterm{def:g11:work-of-force:work}{joule}, zoals elke \omterm{def:g10:energy-conservation:energy}{energie} (\cref{ch:g10:energy-conservation}).
\end{definition}

\begin{example}[Ordegrootten]\label{ex:g11:mechanical-energy:orders}
\begin{itemize}
  \item een voetganger, \qty{70}{kg} met \qty{1.4}{m/s}:
        $E_k = \tfrac12 \times 70 \times 1.4^2 \approx \qty{69}{J}$;
  \item een auto, \qty{1300}{kg} met \qty{130}{km/h} $= \qty{36.1}{m/s}$:
        $E_k \approx \qty{8.5e5}{J}$;
  \item een hogesnelheidstrein, $m = \qty{3.85e5}{kg}$, met
        \qty{320}{km/h} $= \qty{88.9}{m/s}$: $E_k \approx \qty{1.5e9}{J}$;
  \item een geweerkogel, \qty{8.0}{g} met \qty{800}{m/s}:
        $E_k \approx \qty{2.6}{kJ}$.
\end{itemize}
Zeven \omterm{def:g10:orders-of-magnitude:oom}{ordegrootten} scheiden de wandeling van de trein. En $E_k$ groeit
met het \emph{kwadraat} van de snelheid: bij \qty{130}{km/h} draagt de
auto $4$ keer de \omterm{def:g10:energy-conservation:energy}{energie} die hij bij \qty{65}{km/h} heeft --- vier keer
zoveel als zijn remmen moeten wegnemen.
\end{example}

\section{De kinetische-energiestelling}

\begin{theorem}[Kinetische-energiestelling]\label{thm:g11:mechanical-energy:ketheorem}
\index{kinetische-energiestelling}
Voor een lichaam met massa $m$ dat langs een rechte lijn van $A$ naar $B$
beweegt, is de verandering van de \omterm{def:g11:mechanical-energy:kinetic}{kinetische energie} gelijk aan de totale
ontvangen \omterm{def:g11:work-of-force:work}{arbeid} (\cref{ch:g11:work-of-force}):
\[
\Delta E_k = E_k(B) - E_k(A) = \sum W_{AB}(\vect F),
\]
waarbij de som over alle \omterm{def:g10:inertia:force}{krachten} op het lichaam loopt. De stelling
strekt zich uit tot bewegingen langs elke weg, wat we hier eveneens
aannemen.
\end{theorem}
\admitted

\begin{example}[Vrije val, beeldje voor beeldje getoetst]\label{ex:g11:mechanical-energy:freefall}
Een bal van \qty{0.200}{kg} wordt losgelaten en gefilmd; de beeldjes
geven:
\begin{center}\small
\begin{tabular}{lcccc}
$t$ (\unit{s}) & 0.10 & 0.20 & 0.30 & 0.40\\
gevallen hoogte $h$ (\unit{m}) & 0.049 & 0.196 & 0.441 & 0.785\\
snelheid $v$ (\unit{m/s}) & 0.98 & 1.96 & 2.94 & 3.92\\
$E_k = \tfrac12 mv^2$ (\unit{J}) & 0.096 & 0.384 & 0.864 & 1.54\\
$W = mgh$ (\unit{J}) & 0.096 & 0.385 & 0.865 & 1.54
\end{tabular}
\end{center}
Vanaf $E_k = 0$ komt $E_k$ bij elk beeldje overeen met de \omterm{def:g11:work-of-force:work}{arbeid} $mgh$
van het \omterm{def:g10:universal-gravitation:weight}{gewicht}: $\Delta E_k = W(\vect P)$, dus $v^2 = 2gh$ --- de
stelling, geverifieerd tot op de nauwkeurigheid van de aflezingen.
\end{example}

\begin{remark}[Waar het bewijs staat]\label{rem:g11:mechanical-energy:proof}
De algemene uitspraak --- geldig langs elke weg, niet alleen een rechte
lijn --- volgt uit de tweede wet van Newton en een beetje analyse, allebei
verderop in dit boek (\cref{ch:g12:newtons-laws,%
ch:g12:work-and-energy}); tot dan is de toets aan de vrije val ons
onderpand, en nemen we de versie voor gekromde wegen aan en gebruiken we
haar vrijelijk.
\end{remark}

\section{Potentiële energie in het zwaarteveld}

\begin{definition}[Potentiële energie in het zwaarteveld]\label{def:g11:mechanical-energy:potential}
Een lichaam met massa $m$ op hoogte $z$, naar boven gemeten vanaf een
gekozen \emph{referentieniveau}\index{referentieniveau} waar $z = 0$,
slaat de \emph{potentiële energie}\index{potentiële energie}
\[
E_p = m g z, \qquad g = \qty{9.81}{N/kg}
\]
op.
\end{definition}

\begin{remark}[Alleen verschillen tellen]\label{rem:g11:mechanical-energy:reference}
Van $A$ naar $B$ gaan verandert $E_p$ met
$mg(z_B - z_A) = -W_{AB}(\vect P)$: het tegengestelde van de \omterm{def:g11:work-of-force:work}{arbeid} van
het \omterm{def:g10:universal-gravitation:weight}{gewicht} (\cref{ch:g11:work-of-force}). Het \omterm{def:g11:mechanical-energy:potential}{referentieniveau} kies je
vrij --- vloer, zeeniveau, tafelblad --- en het verschuift $E_p$ overal
met dezelfde constante, die uit elk verschil wegvalt. Leg het waar de
getallen het eenvoudigst zijn, en verplaats het niet halverwege een
opgave.
\end{remark}

\begin{example}[Een boek op een plank]\label{ex:g11:mechanical-energy:shelf}
Een boek van \qty{1.2}{kg} ligt \qty{1.8}{m} boven de vloer, die zelf
\qty{9.0}{m} boven de straat ligt. Met de vloer als referentie:
$E_p = 1.2 \times 9.81 \times 1.8 \approx \qty{21}{J}$; met de straat als
referentie: $E_p \approx \qty{127}{J}$. Op de vloer vallen maakt in beide
boekhoudingen \qty{21}{J} vrij.
\end{example}

\section{Mechanische energie en het behoud ervan}

\begin{definition}[Mechanische energie]\label{def:g11:mechanical-energy:em}
De \emph{mechanische energie}\index{mechanische energie} van een lichaam
is
\[
E_m = E_k + E_p = \tfrac12 m v^2 + m g z .
\]
\end{definition}

\begin{theorem}[Behoud van mechanische energie]\label{thm:g11:mechanical-energy:conservation}
\index{behoud van mechanische energie}
Verricht alleen het \omterm{def:g10:universal-gravitation:weight}{gewicht} \omterm{def:g11:work-of-force:work}{arbeid} op een lichaam (geen \omterm{def:g10:inertia:inventory}{wrijving}; andere
\omterm{def:g10:inertia:force}{krachten}, zoals de reactie van een wrijvingsloze \omterm{def:g10:relative-motion:trajectory}{baan}, loodrecht op de
beweging), dan blijft $E_m$ behouden: voor elke twee punten $A$ en $B$
van de beweging geldt
\[
\tfrac12 m v_A^2 + m g z_A = \tfrac12 m v_B^2 + m g z_B .
\]
\end{theorem}

\begin{proof}
Volgens de kinetische-energiestelling is
$E_k(B) - E_k(A) = W_{AB}(\vect P) = mg(z_A - z_B) = E_p(A) - E_p(B)$:
wat $E_k$ wint, verliest $E_p$.
\end{proof}

\begin{example}[Waterglijbaan]\label{ex:g11:mechanical-energy:slide}
Een kind vertrekt uit rust boven aan een wrijvingsloze glijbaan met
hoogte $h = \qty{3.2}{m}$. Met de referentie onderaan geldt
$mgh = \tfrac12 mv^2$, dus $v = \sqrt{2gh} \approx \qty{7.9}{m/s}$
(\qty{29}{km/h}) --- wat de massa ook is \emph{en welke vorm de \omterm{def:g10:relative-motion:trajectory}{baan} ook
heeft}: recht, gebogen of spiraalvormig, alleen de daling telt mee.
\end{example}

\begin{example}[Slinger]\label{ex:g11:mechanical-energy:pendulum}
Een slingerbol wordt opzij getrokken tot hij $h = \qty{12}{cm}$ boven
zijn laagste punt staat, en dan losgelaten. In de uiterste standen is
$v = 0$; onderaan is de hele $mgh$ kinetisch,
$v = \sqrt{2gh} \approx \qty{1.5}{m/s}$; aan de andere kant klimt de bol
weer tot precies \qty{12}{cm}.
\end{example}

\begin{omfigure}
\begin{tikzpicture}[scale=0.95, font=\small]
  \coordinate (O) at (0,0); \draw (-0.5,0) -- (0.5,0);
  \fill (O) circle (1.5pt);
  \draw[dashed, omExo] (230:2.2) arc (230:310:2.2);
  \draw (O) -- (230:2.2); \draw (O) -- (270:2.2); \draw (O) -- (310:2.2);
  \fill[omDef] (230:2.2) circle (2.5pt); \fill[omDef] (270:2.2)
    circle (2.5pt); \fill[omDef] (310:2.2) circle (2.5pt);
  \draw[-Stealth, omThm, thick] (270:2.2) -- ++(1.15,0)
    node[below] {$v_{\max}$};
  \node[above left] at (230:2.2) {$v = 0$};
  \node[above right] at (310:2.2) {$v = 0$};
  \draw[dashed, omExo] (230:2.2) -- ++(-0.9,0)
    (270:2.2) -- ++(-2.4,0);
  \draw[Stealth-Stealth] (-2.2,-2.2) -- (-2.2,-1.685)
    node[midway, left] {$h$};
\end{tikzpicture}

{\small De slinger ruilt $E_p$ tegen $E_k$ en terug: in de uiterste
standen alles potentieel, onderaan alles kinetisch, en twee keer
dezelfde stijging $h$.}
\end{omfigure}

\begin{omfigure}
\begin{tikzpicture}[scale=0.75, font=\small]
  \draw[thick] plot [smooth, tension=0.7] coordinates
    {(0,3.2) (1.2,1.6) (2.4,0.5) (3.6,1.0) (4.6,2.0) (5.6,0.9)};
  \fill (0,3.2) circle (2pt) node[above right] {$A$};
  \fill (2.4,0.5) circle (2pt) node[below] {$B$}
    (4.6,2.0) circle (2pt) node[above] {$C$};
  \draw[omExo] (-0.2,0) -- (6.0,0);
  \fill[omProp!75] (7.0,0) rectangle (7.45,1.6);
  \fill[omProp!75] (7.9,0) rectangle (8.35,0.25)
    (8.8,0) rectangle (9.25,1.0);
  \fill[omDef!75] (7.9,0.25) rectangle (8.35,1.6)
    (8.8,1.0) rectangle (9.25,1.6);
  \draw[omExo] (6.7,0) -- (9.6,0);
  \draw[dashed, omThm] (6.7,1.6) -- (9.6,1.6) node[above left] {$E_m$};
  \node[below] at (7.225,0) {$A$}; \node[below] at (8.125,0) {$B$};
  \node[below] at (9.025,0) {$C$};
  \fill[omProp!75] (6.9,3.0) rectangle +(0.24,0.24);
  \fill[omDef!75] (6.9,2.5) rectangle +(0.24,0.24);
  \node[right] at (7.2,3.12) {$E_p$}; \node[right] at (7.2,2.62) {$E_k$};
\end{tikzpicture}

{\small Staafdiagram van de \omterm{def:g10:energy-conservation:energy}{energie} van een wrijvingsloze achtbaan: in
elke stand stapelen de staven $E_p$ en $E_k$ tot hetzelfde totaal $E_m$
(gestippelde lijn).}
\end{omfigure}

\begin{example}[Landingssnelheid bij het schansspringen]\label{ex:g11:mechanical-energy:skijump}
Een schansspringer verlaat de schans met $v_A = \qty{25}{m/s}$ en landt
$\Delta z = \qty{40}{m}$ lager. Zonder luchtweerstand geldt
$v_B = \sqrt{v_A^2 + 2g\,\Delta z} = \sqrt{625 + 785} \approx
\qty{38}{m/s}$ --- zonder dat je iets van de vluchtbaan hoeft te weten.
\end{example}

\begin{omfigure}
\begin{tikzpicture}[scale=0.75, font=\small]
  \draw[thick] (0,3.8) .. controls (1.2,2.7) and (1.9,2.4) .. (3.1,2.3);
  \draw[thick] (3.1,1.95) .. controls (4.2,1.55) and (5.5,1.0) .. (7.2,0);
  \draw[dashed, omThm, thick] (3.1,2.3) .. controls (4.3,2.25) and
    (5.2,1.7) .. (5.9,0.75);
  \fill (3.1,2.3) circle (2pt) node[above] {$A$}
    (5.9,0.75) circle (2pt) node[above right] {$B$};
  \draw[dashed, omExo] (3.1,2.3) -- (8.2,2.3) (5.9,0.75) -- (8.2,0.75);
  \draw[Stealth-Stealth] (8.0,0.75) -- (8.0,2.3)
    node[midway, right] {$\Delta z$};
\end{tikzpicture}

{\small Schansspringen: welke kromme er ook van $A$ naar $B$ wordt
gevlogen, alleen de daling $\Delta z$ komt in de energiebalans voor.}
\end{omfigure}

\section{Dissipatie door wrijving}

\begin{definition}[Dissipatie]\label{def:g11:mechanical-energy:dissipation}
\omterm{def:g10:energy-conservation:energy}{Energie} die de mechanische rekening verlaat, heet
\emph{gedissipeerd}\index{dissipatie}: ze duikt weer op als \omterm{def:g10:energy-conservation:forms}{thermische
energie} in de wrijvende oppervlakken en in de lucht. De totale \omterm{def:g10:energy-conservation:energy}{energie}
blijft behouden (\cref{ch:g10:energy-conservation}); de \omterm{def:g11:mechanical-energy:em}{mechanische
energie} alleen niet.
\end{definition}

\begin{proposition}[Het tekort meet de wrijving]\label{prop:g11:mechanical-energy:deficit}
Verricht ook een wrijvingskracht $\vect f$ \omterm{def:g11:work-of-force:work}{arbeid} langs de beweging, dan
geldt
\[
\Delta E_m = E_m(B) - E_m(A) = W_{AB}(\vect f) < 0:
\]
het verlies aan \omterm{def:g11:mechanical-energy:em}{mechanische energie} is gelijk aan de \omterm{def:g11:work-of-force:work}{arbeid} van de
\omterm{def:g10:inertia:inventory}{wrijving}.
\end{proposition}

\begin{proof}
De kinetische-energiestelling luidt nu
$\Delta E_k = W_{AB}(\vect P) + W_{AB}(\vect f) = -\Delta E_p +
W_{AB}(\vect f)$; breng $\Delta E_p$ naar de andere kant. De \omterm{def:g10:inertia:inventory}{wrijving}
werkt de beweging tegen, dus is haar \omterm{def:g11:work-of-force:work}{arbeid} negatief.
\end{proof}

\begin{example}[Doorlichting van een slee]\label{ex:g11:mechanical-energy:toboggan}
Een kind met slee ($m = \qty{40}{kg}$) vertrekt uit rust, daalt
$h = \qty{5.0}{m}$ over \qty{12}{m} sneeuw, en komt aan met
\qty{7.0}{m/s} in plaats van de wrijvingsloze
$\sqrt{2gh} \approx \qty{9.9}{m/s}$. Met de referentie onderaan:
$\Delta E_m = \tfrac12 \times 40 \times 7.0^2 - 40 \times 9.81
\times 5.0 = 980 - 1962 = \qty{-982}{J}$. Dus is
$W(\vect f) = \qty{-982}{J}$ over \qty{12}{m}: een gemiddelde
wrijvingskracht $f = 982/12 \approx \qty{82}{N}$, en \qty{982}{J} aan
warmte in de glijders en de sneeuw.
\end{example}

\begin{omfigure}
\begin{tikzpicture}
\begin{axis}[omaxis, width=8.6cm, height=4.6cm,
    xlabel={$t$ (\unit{s})}, ylabel={$E_m$ (\unit{J})},
    xmin=0, xmax=10, ymin=0, ymax=1250, legend pos=south west]
  \addplot[omThm, thick, domain=0:10, samples=2] {1000};
  \addplot[omDef, thick, domain=0:10, samples=60] {1000*exp(-0.18*x)};
  \legend{zonder wrijving, met wrijving}
\end{axis}
\end{tikzpicture}

{\small $E_m$ tijdens een beweging: vlak zonder \omterm{def:g10:inertia:inventory}{wrijving}, dalend met
\omterm{def:g10:inertia:inventory}{wrijving} --- de daling op elk ogenblik is de tot dan toe geproduceerde
warmte.}
\end{omfigure}

\begin{method}[Energieboekhouding]\label{met:g11:mechanical-energy:bookkeeping}
\begin{enumerate}
  \item Kies een \omterm{def:g11:mechanical-energy:potential}{referentieniveau} en houd het aan; benoem $A$ (gegevens)
        en $B$ (vraag).
  \item Schrijf $E_m = \tfrac12 mv^2 + mgz$ op in $A$ en in $B$.
  \item Geen \omterm{def:g10:inertia:inventory}{wrijving}: $E_m(A) = E_m(B)$; los op --- doorgaans
        $v_B = \sqrt{v_A^2 + 2g(z_A - z_B)}$.
  \item Wel \omterm{def:g10:inertia:inventory}{wrijving}: $E_m(B) - E_m(A) = -f\ell$ ($\ell$ = weglengte);
        los de onbekende op.
  \item Controleer: de staven $E_k$ en $E_p$ stapelen tot $E_m$, dat vlak
        blijft tenzij er \omterm{def:g10:inertia:inventory}{wrijving} is.
\end{enumerate}
\end{method}

\section{Oefeningen}

\begin{exercise}[$\star$]\label{exo:g11:mechanical-energy:1}
Bereken de \omterm{def:g11:mechanical-energy:kinetic}{kinetische energie} van (a) een scooter met berijder van
\qty{90}{kg} bij \qty{25}{km/h}; (b) een auto van \qty{1200}{kg} bij
\qty{130}{km/h}. (c) Met welke factor verandert $E_k$ als de snelheid
verdubbelt? En als ze verdrievoudigt?
\end{exercise}

\begin{exercise}[$\star$]\label{exo:g11:mechanical-energy:2}
Een klimster van \qty{75}{kg} wint \qty{850}{m} hoogte. Bereken de winst
aan \omterm{def:g11:mechanical-energy:potential}{potentiële energie}. Hangt die van het \omterm{def:g11:mechanical-energy:potential}{referentieniveau} af? En van de
route?
\end{exercise}

\begin{exercise}[$\star$]\label{exo:g11:mechanical-energy:3}
Een bloempot valt \qty{20}{m} van een balkon. Bepaal, met verwaarloosbare
luchtweerstand, zijn snelheid bij de grond in \unit{m/s} en in
\unit{km/h}. Waar is de massa gebleven?
\end{exercise}

\begin{exercise}[$\star$]\label{exo:g11:mechanical-energy:4}
Een auto van \qty{1200}{kg} remt in \qty{50}{m} van \qty{25}{m/s} tot
stilstand. Bepaal met de kinetische-energiestelling de \omterm{def:g11:work-of-force:work}{arbeid} van de
remmen, en daarna de gemiddelde remkracht.
\end{exercise}

\begin{exercise}[$\star$]\label{exo:g11:mechanical-energy:5}
Een slingerbol wordt \qty{8.0}{cm} boven zijn laagste punt losgelaten.
Bepaal zijn snelheid onderaan en de hoogte die hij aan de andere kant
bereikt.
\end{exercise}

\begin{exercise}[$\star\star$]\label{exo:g11:mechanical-energy:6}
Met \cref{ex:g11:mechanical-energy:orders}: met welke factor overtreft de
\omterm{def:g11:mechanical-energy:kinetic}{kinetische energie} van de trein die van de kogel? Tot welke hoogte zou
zijn eigen $E_k$ de trein kunnen tillen ($h = v^2/2g$)?
\end{exercise}

\begin{exercise}[$\star\star$]\label{exo:g11:mechanical-energy:7}
Twee wrijvingsloze glijbanen dalen van hetzelfde platform van
\qty{3.0}{m}, de ene recht en steil, de andere lang en kronkelend.
Vergelijk de aankomstsnelheden (bereken ze) en de aankomsttijden (zonder
rekenwerk).
\end{exercise}

\begin{exercise}[$\star\star$]\label{exo:g11:mechanical-energy:8}
Een bal wordt recht omhoog gegooid met \qty{12}{m/s}. Bepaal de maximale
hoogte, en daarna de snelheid op \qty{4.0}{m} --- op de weg omhoog en op
de weg omlaag.
\end{exercise}

\begin{exercise}[$\star\star$]\label{exo:g11:mechanical-energy:9}
Een schansspringer verlaat de schans met \qty{23}{m/s} en landt
\qty{35}{m} lager. Voorspel de landingssnelheid zonder luchtweerstand; is
de werkelijke waarde groter of kleiner, en waarom?
\end{exercise}

\begin{exercise}[$\star\star$]\label{exo:g11:mechanical-energy:10}
Een slee van \qty{45}{kg} vertrekt uit rust, daalt \qty{6.0}{m} over
\qty{15}{m} helling en komt aan met \qty{8.0}{m/s}. Bepaal de verloren
\omterm{def:g11:mechanical-energy:em}{mechanische energie} en daarna de gemiddelde wrijvingskracht. Waar is de
\omterm{def:g10:energy-conservation:energy}{energie} gebleven?
\end{exercise}

\begin{exercise}[$\star\star$]\label{exo:g11:mechanical-energy:11}
De pen van Galilei: een slinger die \qty{15}{cm} boven zijn laagste punt
wordt losgelaten, treft in de verticale stand een pen die de bovenste
helft van het touw blokkeert. Bepaal de snelheid onderaan en de stijging
voorbij de pen; wat verandert de pen, en wat niet?
\end{exercise}

\begin{exercise}[$\star\star\star$]\label{exo:g11:mechanical-energy:12}
Een skateboarder laat zich van \qty{2.5}{m} in een halfpipe zakken; elke
oversteek (van kant tot kant) dissipeert $10\%$ van de \omterm{def:g11:mechanical-energy:em}{mechanische
energie}. Bepaal de snelheid onderaan de eerste afdaling, en daarna ---
door de opeenvolgende stijghoogten te berekenen --- de eerste oversteek
die onder \qty{1.0}{m} eindigt.
\end{exercise}

\begin{exercise}[$\star\star\star$]\label{exo:g11:mechanical-energy:13}
Vanaf de rand van een klif van \qty{45}{m} wordt een steen met
\qty{15}{m/s} weggegooid --- omhoog, horizontaal of onder welke hoek dan
ook. Toon aan dat de landingssnelheid in alle gevallen dezelfde is en
bereken haar. Wat hangt er \emph{wel} van de hoek af?
\end{exercise}

\begin{exercise}[$\star\star\star$]\label{exo:g11:mechanical-energy:14}
Een afdalingsskiër van \qty{80}{kg} vertrekt uit rust en daalt
\qty{120}{m} over een piste van \qty{800}{m}, tegen een constante
wrijvingskracht van \qty{65}{N} in. Bepaal de aankomstsnelheid en het
deel van $E_m$ dat \omterm{def:g11:mechanical-energy:dissipation}{gedissipeerd} is.
\end{exercise}

\begin{exercise}[$\star\star\star$]\label{exo:g11:mechanical-energy:15}
Een pompaccumulatiecentrale tilt $V = \qty{2.0e6}{m^3}$ water
(\qty{1000}{kg} per \unit{m^3}) \qty{300}{m} omhoog. Bereken de
opgeslagen \omterm{def:g10:energy-conservation:energy}{energie} in \omterm{def:g11:work-of-force:work}{joule} en in \unit{kW.h}; als die met een \omterm{def:g10:energy-conservation:efficiency}{rendement}
van $85\%$ weer wordt afgeturbineerd
(\cref{ch:g10:energy-conservation}), hoeveel huishoudens die
\qty{10}{kW.h} per dag gebruiken, voedt ze dan één dag lang?
\end{exercise}

\section{Probleem: De looping ontwerpen}

\begin{problem}[{Weekendopgave --- een energiedoorlichting van een
achtbaan: de starthoogte die een looping voorschrijft, de snelheden die
ze belooft, de wrijving die twee sensoren opbiechten, en de remmen die de
rekening sluiten --- de hele rit hangend aan één getal}]
\label{pb:g11:mechanical-energy:1}
Een trein met massa $m = \qty{2.0e3}{kg}$ wordt uit rust losgelaten van
een startheuvel met hoogte $H$ en rolt, zonder motor, door een verticale
cirkelvormige looping met straal $R = \qty{8.0}{m}$ waarvan de top op
hoogte $2R$ ligt. De analyse van de cirkelbeweging (volgend jaar) levert
het ene gegeven dat we lenen: bovenin de looping drukt de trein alleen op
de rails als $v_{\text{top}}^2 \geq gR$. Deel I en II verwaarlozen de
\omterm{def:g10:inertia:inventory}{wrijving}.

\medskip
\noindent\textbf{Deel I --- De heuvel die de looping voorschrijft.}
\begin{enumerate}
  \item Bereken de minimaal veilige snelheid boven in de looping.
  \item Met de grond als \omterm{def:g11:mechanical-energy:potential}{referentieniveau}: bereken de \omterm{def:g11:mechanical-energy:em}{mechanische energie}
        van de trein boven in de looping bij die minimale snelheid.
  \item Leid $H_{\min}$ af en toon aan dat $H_{\min} = 5R/2$ --- zonder
        massa en zonder $g$.
  \item Het team kiest $H = \qty{24}{m}$: bereken de werkelijke snelheid
        boven in de looping en de verhouding $v_{\text{top}}^2/(gR)$.
  \item Waar is $m$ gebleven? Waarom past het ontwerp even goed bij volle
        als bij lege treinen?
\end{enumerate}

\medskip
\noindent\textbf{Deel II --- Wat de hoogte belooft} ($H = \qty{24}{m}$).
\begin{enumerate}[resume]
  \item Wat is de snelheid onder in de looping ($z = 0$), in \unit{m/s} en
        \unit{km/h}?
  \item Wat is de snelheid halverwege de looping ($z = R$)?
  \item Bereken $E_p$ en $E_k$ bij de start, onderaan, halverwege en
        boven in de looping; ga na dat elk paar tot $mgH$ optelt.
  \item Na de looping volgt een bult van \qty{12}{m}: wat is de snelheid
        op de top ervan?
  \item Geef de algemene formule voor $v(z)$; waar is de rit het snelst,
        en waar het traagst?
\end{enumerate}

\medskip
\noindent\textbf{Deel III --- De doorlichting van de \omterm{def:g10:inertia:inventory}{wrijving}.}
Twee sensoren op dezelfde hoogte $z = \qty{2.0}{m}$, met \qty{60}{m}
rails ertussen, meten $v_1 = \qty{20.4}{m/s}$ en daarna
$v_2 = \qty{19.2}{m/s}$.
\begin{enumerate}[resume]
  \item Bereken de verandering van de \omterm{def:g11:mechanical-energy:em}{mechanische energie} tussen de
        sensoren. Wat maakte de opstelling op gelijke hoogte zo slim?
  \item Leid de \omterm{def:g11:work-of-force:work}{arbeid} van de \omterm{def:g10:inertia:inventory}{wrijving} en de gemiddelde wrijvingskracht
        af.
  \item Waar is de ontbrekende \omterm{def:g10:energy-conservation:energy}{energie} gebleven? Noem de vorm en de
        plaatsen.
  \item Welk deel van de $E_m$ van de trein gaat er per \qty{60}{m}
        verloren? Waarom zou de kale $H_{\min}$ in de echte, wrijvende
        wereld onveilig zijn?
  \item Tussen start en looptop ligt ongeveer \qty{90}{m} rails: schat
        $E_m$ daar; redt de marge uit vraag 4 de looping nog?
\end{enumerate}

\medskip
\noindent\textbf{Deel IV --- De rekening sluiten.}
De laatste sensor, bij de ingang van het rechte horizontale remtraject
($z = 0$), leest \qty{18.0}{m/s}.
\begin{enumerate}[resume]
  \item Bereken de \omterm{def:g11:mechanical-energy:kinetic}{kinetische energie} van de trein daar.
  \item De remmen moeten hem in $d = \qty{45}{m}$ tot stilstand brengen:
        bereken de vereiste (constante) remkracht.
  \item Comfortcontrole: bereken de vertraging $a = F/m$ en vergelijk met
        $g$.
  \item Hoeveel \omterm{def:g10:energy-conservation:energy}{energie} heeft de \omterm{def:g10:inertia:inventory}{wrijving} met de rails over de hele rit
        \omterm{def:g11:mechanical-energy:dissipation}{gedissipeerd}? Ga na dat \omterm{def:g10:inertia:inventory}{wrijving} $+$ remmen $= mgH$, tot op de
        \omterm{def:g11:work-of-force:work}{joule}.
  \item De clou, in één zin: welk enkel getal legde de veiligheid van de
        looping, elke snelheid, de wrijvingsrekening en de remkracht
        vast --- en wat was de prijs van het negeren van de \omterm{def:g10:inertia:inventory}{wrijving}?
\end{enumerate}
\end{problem}
